\chapter{Double Integrals in Polar Form}
\label{chap:double-integrals-in-polar-form}

\section{Double Integrals}
\label{sec:double-integrals-polar-form}

\begin{definition}[Double Integral in Polar Coordinates]
  \label{def:polar-double-integral}
  If \(f(x,y)\) is continuous on a region \(R\), and we convert to polar coordinates using:
  \[
    x = r\cos\theta, \quad y = r\sin\theta
  \]
  then the double integral becomes:
  \begin{equation}\label{eq:polar-double-integral}
    \iint_R f(x,y)\,dx\,dy = \iint_R f(r\cos\theta, r\sin\theta)\,r\,dr\,d\theta
  \end{equation}
\end{definition}

\begin{remark}[Why \(r\) Appears]
  The extra factor \(r\) accounts for the Jacobian determinant when changing from rectangular to polar coordinates.
\end{remark}

\section{Relating the Double Integrals in Cartesian and Polar Forms}
\label{sec:relating-the-double-integrals-in-cartesian-and-polar-forms}

\begin{equation}\label{eq:jacobian-cartesian-polar}
  dx\,dy = r\,dr\,d\theta
\end{equation}

\begin{explanation}[Jacobian of the Transformation]
  The Jacobian determinant for the transformation from \((x,y)\) to \((r,\theta)\) is:
  \[
    J = \left| \frac{\partial(x,y)}{\partial(r,\theta)} \right| = 
    \begin{vmatrix}
      \cos\theta & -r\sin\theta \\
      \sin\theta & r\cos\theta
    \end{vmatrix}
    = r
  \]
  Hence, \(dx\,dy = r\,dr\,d\theta\).
\end{explanation}

\section{Converting Double Integrals from Cartesian to Polar Forms}
\label{sec:converting-double-integrals-from-cartesian-to-polar-forms}

\begin{exercise}
  Convert the integral
  \[
    \int_{-1}^{1} \int_{0}^{\sqrt{1-x^2}} 1\,dy\,dx
  \]
  to polar form.
\end{exercise}

\begin{answer}[Converting Circle Integral to Polar Form]
  The given region is the upper half of the unit disk:
  \[
    R = \{ (x,y): -1 \le x \le 1, 0 \le y \le \sqrt{1 - x^2} \}
  \]
  This corresponds to:
  \[
    R = \{ (r,\theta): 0 \le r \le 1,\ 0 \le \theta \le \pi \}
  \]

  Converting the integral:
  \begin{align*}
    \int_{-1}^{1} \int_{0}^{\sqrt{1-x^2}} 1\,dy\,dx 
    &= \iint_R 1\,dx\,dy \\
    &= \int_0^{\pi} \int_0^1 1 \cdot r\,dr\,d\theta \\
    &= \int_0^{\pi} \left[ \frac{1}{2}r^2 \right]_0^1 d\theta \\
    &= \int_0^{\pi} \frac{1}{2}\,d\theta = \left[ \frac{1}{2}\theta \right]_0^{\pi} = \frac{\pi}{2}
  \end{align*}
\end{answer}

\begin{remark}[When to Use Polar Coordinates]
  Use polar form when the region is circular or radial symmetry is present, as integration becomes simpler.
\end{remark}

\begin{proof}[Equivalence of Cartesian and Polar Forms]
  Let \(R\) be a region described easily in polar coordinates. Then:
  \[
    \iint_R f(x,y)\,dx\,dy = \iint_R f(r\cos\theta, r\sin\theta)\cdot r\,dr\,d\theta
  \]
  This is justified by the substitution rule in multiple integrals and the Jacobian determinant of the coordinate transformation.
\end{proof}
